\section{Non-Localization}


This section organizes why consciousness cannot be localized to a specific site.

Non-localization here does not mean consciousness has no neural substrate.

Conscious states clearly depend on bodily and brain activity.

But that dependency does not take the form of consciousness "existing" at a single site.

In the position of this paper, consciousness is not a location but a temporal regime that opens when speed difference, coupling viscosity, internal delay, meaning gradient, and log referencing enter a certain relation.

An attempt to localize consciousness is therefore likely searching not for consciousness itself but for the residue that remains after consciousness has passed through, or the local substrate that participates in consciousness.

\subsection{The Mistaken Search for a Seat}


Consciousness research has a recurring tendency to search for a "seat of consciousness."

Which brain region generates consciousness? Which circuit produces subjective experience? Which signal corresponds to consciousness itself?

These questions seem natural at first glance.

But behind them lies the generation model: consciousness is produced somewhere, and what is produced exists somewhere.

This paper does not adopt that presupposition.

Consciousness is not a product but an opening state.

It is the state in which referenceability has opened between processing and logs, and its establishment involves multiple time scales, plasticity gate, meaning gradient, and non-closure loop.

The question of placing consciousness at a single point therefore has the wrong object from the start.

\subsection{Local Substrates and Non-Local Regimes}


However, the non-localization of consciousness does not mean consciousness has no local substrate.

Fast-layer processing has corresponding neural activity.

Slow-layer log sedimentation also has corresponding neural and bodily substrate.

Reporting pathways, memory pathways, sensory processing, motor preparation, and verbalization each have their own local substrates.

These are important.

But none of them individually is consciousness itself.

Consciousness is the regime that opens when local substrates temporally couple and enter a log-referenceable relation.

What is required in consciousness research is therefore not the denial of local substrates.

What is required is not conflating local substrates with the non-local regime.

Brain regions are involved.

But the conscious state obtains not as regions but as the relation among regions, temporal layers, and logs.

\subsection{Why Localization Converts Process into Residue}


Localization involves the operation of fixing an object.

To assign a phenomenon to a specific location, the phenomenon must be treated as a stable object.

But consciousness is a live opening state, not a completely fixed object.

As seen in §13, measurement does not preserve qualia or the conscious state as they are; it sediments them into the slow layer and converts them into residue.

Localization is the same.

What localization yields is not live consciousness itself but the residue, trace, correlation, or reportable log that was involved in the conscious state.

The presence of activity in a certain region is therefore important.

But it does not mean consciousness exists in that region.

It shows that the region is participating as part of the log-referencing loop or its residue.

Localization is not meaningless.

But what it captures is not consciousness itself but the local conditions that were involved when consciousness opened.

\subsection{Consciousness as Temporal Relation}


In the framework of this paper, consciousness is not a spatial location but a temporal relation.

The fast layer anticipates the future. The slow layer retains the past. The plasticity gate regulates what sediments and what flows away between them. The meaning gradient directs which logs are easily referenced.

When this relation enters a certain range, the conscious state opens.

The appropriate form of the question about consciousness is therefore not "where is it?" but "under what temporal conditions does processing become log-referenceable?"

This condition was expressed as the consciousness window:

\begin{equation}
0 < |\Delta v| < \Delta v_{\max}, \quad \mu_{\min} < \mu < \mu_{\max}, \quad \tau_{\min} < \tau < \tau_{\max}, \quad \mathcal{T}_{\min} < \mathcal{T}_{\Phi} < \mathcal{T}_{\max}
\end{equation}

This window is not a point in space.

It is a dynamic region in which multiple processing layers and log structures are temporally coupled.

\subsection{Neural Correlates Are Not Consciousness Itself}


Searching for the neural correlates of consciousness is important.

But neural correlates are not consciousness itself.

Neural correlates indicate the conditions under which conscious states open, or the traces left after conscious states have opened.

For example, if a certain neural activity corresponds to a particular conscious report, that correspondence is important information.

But this alone does not allow saying that neural activity is consciousness itself.

The reason is that the same activity can carry different meaning depending on reporting conditions, attentional state, meaning gradient, delay, and plasticity gate state.

This paper does not deny neural correlates.

Rather, it reads neural correlates as part of the log-referencing loop.

Neural activity is not the sole seat that generates consciousness but one element of the temporal non-closure structure through which consciousness opens.

\subsection{Distributed Does Not Mean Diffuse}


Saying that consciousness is non-local can be misunderstood as implying that consciousness is spread vaguely throughout the brain.

But the position of this paper is not such a diffusion model.

Non-local does not mean locationless; it means not reducible to a single location.

The conscious state obtains when multiple local processes enter a specific temporal relation.

Consciousness is therefore not everywhere.

Rather, it opens only under quite strict conditions.

Speed difference is required. Viscosity is required. Delay is required. Meaning gradient is required. Plasticity gate is required. And these must reach neither complete closure nor complete separation.

In this sense, consciousness is not diffuse but conditioned.

\subsection{From Localization to Conditions}


There is no need to abandon the localization question entirely.

But it is not the final form of the question.

The question "which region is active?" should be transformed into "which local substrates are participating in the log-referencing loop?"

This question is then extended as follows: which time scales are involved; which speed differences are maintained; how much delay is there; which logs have become referenceable; which meaning gradients are directing referencing; which plasticity gates are opening and closing.

Changing the questions in this way transforms consciousness from a "location-finding" problem into a condition-arrangement problem.

This is the shift in questioning that this paper proposes.

\subsection{Summary}


This section organized the non-localization of consciousness.

\begin{itemize}
\item Consciousness has neural substrates.
\item But the conscious state itself does not localize to a single site.
\item What localization captures is not live consciousness but local substrates or sedimented residue.
\item Consciousness is the temporal relation formed by speed difference, coupling viscosity, internal delay, meaning gradient, and plasticity gate.
\item Neural correlates are important, but they should be interpreted not as consciousness itself but as part of the log-referencing loop.
\item Non-local does not mean diffuse; it means not reducible to a single site.
\item The appropriate question is not "where is consciousness?" but "under what conditions does the conscious state open?"
\end{itemize}

The non-localization in this paper is therefore not mystification.

It is the consequence of treating consciousness not as a spatial object but as a temporal non-closure regime.
